\documentclass[a4paper, serif, 10pt]{moderncv}

%% moderncv
% color and style
\moderncvcolor{black}
\moderncvstyle{banking}

%% page layout/margins
\usepackage[top=2.5cm, bottom=2.5cm, left=1.5cm, right=1.5cm]{geometry}

\nopagenumbers{}

%% Babel config for Italian language
% ref:
% - https://latex3.github.io/babel/guides/locale-italian.html
\usepackage[italian]{babel}

%% fontspec
\usepackage{fontspec}

\IfFontExistsTF{Linux Libertine}{
  \setmainfont[Ligatures={TeX, Common}, Numbers=OldStyle]{Linux Libertine}
}{}
\IfFontExistsTF{Linux Libertine O}{
  \setmainfont[Ligatures={TeX, Common}, Numbers=OldStyle]{Linux Libertine O}
}{}

%% CTeX
% CJK support, used to show CJK characters in the resume
%
% - heading=false: disable ctex heading
% - punct=kaiming: use kaiming punctuations styles for CJK
% - scheme=plain: use plain scheme, do not override \normalsize font size
% - space=auto: space settings for CJK characters
%
% CJK fonts are resolved per locale language, e.g. Noto Serif/Sans CJK JP
% for Japanese resumes, with fallback to the next region if not installed
%
% ref:
% - http://ctan.mirrorcatalogs.com/language/chinese/ctex/ctex.pdf
\usepackage[UTF8, heading=false, punct=kaiming, scheme=plain, space=auto]{ctex}

\IfFontExistsTF{Noto Serif CJK SC}{
  \setCJKmainfont{Noto Serif CJK SC}
}{}
\IfFontExistsTF{Noto Sans CJK SC}{
  \setCJKsansfont{Noto Sans CJK SC}
}{}

%% line spacing
\usepackage{setspace}
\setstretch{1.125}

%% URL styling - use normal text instead of monospace
\urlstyle{same}

%% Patch \httplink and \httpslink to handle full URLs
% The original moderncv macros always prepend http:// or https:// to URLs.
% This causes issues when the URL already has a protocol (e.g., https://example.com)
% resulting in malformed URLs like https://https://example.com.
% This patch checks if the URL already contains "://" and uses it directly if so.
% Note: We use \str_set:Nx (x-expansion) to expand macros like \@homepage first.
\ExplSyntaxOn
\str_new:N \g_yamlresume_url_str

\RenewDocumentCommand{\httpslink}{O{}m}{
  \str_set:Nx \g_yamlresume_url_str {#2}
  \str_if_in:NnTF \g_yamlresume_url_str {://}
    { \url{#2} }
    { \url{https://#2} }
}

\RenewDocumentCommand{\httplink}{O{}m}{
  \str_set:Nx \g_yamlresume_url_str {#2}
  \str_if_in:NnTF \g_yamlresume_url_str {://}
    { \url{#2} }
    { \url{http://#2} }
}
\ExplSyntaxOff

\name{Giulia Bianchi}{}
\title{Customer Success Manager | SaaS e trasformazione digitale}
\phone[mobile]{+39 345 678 9012}
\email{giulia.bianchi.csm@email.it}
\homepage{https://www.giuliabianchi-csm.it}

\address{Via della Moscova 18, Milano, Lombardia, Italia, 20121}{}{}

\extrainfo{{\small \faLinkedin}\ \href{https://www.linkedin.com/in/giulia-bianchi-csm}{@giulia-bianchi-csm} {} {} {} • {} {} {} 
{\small \faMedium}\ \href{https://medium.com/@giuliabianchi\_csm}{@giuliabianchi\_csm}}

\begin{document}

\maketitle

\section{Informazioni personali}

\cvline{}{Customer Success Manager con oltre 8 anni di esperienza nella gestione di portafogli clienti B2B SaaS, con una solida propensione alla fidelizzazione, all'adozione del prodotto e alla crescita del valore del cliente.
%
\begin{itemize}
\item Responsabile di un portafoglio da 2,5 milioni di euro di ARR, con una riduzione del churn del 18\% in due anni.
\item Esperta nella costruzione di processi di onboarding scalabili, piani di successo del cliente e business review trimestrali.
\item Forte orientamento ai dati, alla collaborazione cross-funzionale e alla comunicazione con stakeholder senior.
\item Abituata a lavorare in contesti internazionali e a tradurre le esigenze del cliente in roadmap di prodotto.
\end{itemize}}
\section{Formazione}

\cventry{ott 2012–lug 2014}
        {Laurea magistrale, Marketing e Comunicazione d'Impresa, Punteggio: 110/110 con lode}
        {Università Commerciale Luigi Bocconi}
        {\url{https://www.unibocconi.it}}
        {}
        {\begin{itemize}
\item Percorso specialistico focalizzato sulla gestione della relazione con il cliente e sulla misurazione della soddisfazione.
\item Tesi sperimentale sull'impatto della customer experience sulla fidelizzazione nei servizi digitali.
\item Progetti di gruppo con aziende del settore tecnologico e retail.
\end{itemize}
\textbf{Corsi}: Marketing strategico, Comportamento del consumatore, Ricerche di mercato, Gestione della relazione con il cliente, Business analytics, Comunicazione d'impresa}

\cventry{set 2009–lug 2012}
        {Laurea triennale, Scienze della Comunicazione, Punteggio: 108/110}
        {Università degli Studi di Milano}
        {\url{https://www.unimi.it}}
        {}
        {\begin{itemize}
\item Solida base in comunicazione interpersonale, analisi dei bisogni e tecniche di negoziazione.
\item Attività di laboratorio su customer journey e ricerca qualitativa.
\end{itemize}
\textbf{Corsi}: Sociologia della comunicazione, Psicologia del consumatore, Statistica, Linguaggi digitali, Public speaking}
\section{Esperienza professionale}

\cventry{gen 2021–Presente}
        {Senior Customer Success Manager}
        {NexaCloud Italia}
        {\url{https://www.nexacloud.it}}
        {}
        {\begin{itemize}
\item Gestione di un portafoglio di 45 clienti enterprise e mid-market, per un ARR complessivo di 2,5 milioni di euro.
\item Definizione e implementazione di piani di successo del cliente con obiettivi misurabili su adozione, valore e rinnovo.
\item Riduzione del churn del 18\% e aumento del tasso di rinnovo al 94\% in due anni.
\item Conduzione di business review trimestrali con C-level, HR e responsabili IT per allineare le priorità e dimostrare il ROI.
\item Collaborazione con Product e Sales per raccogliere feedback, identificare opportunità di upselling e migliorare la roadmap.
\item Mentoring di due Customer Success Manager junior e definizione delle best practice di onboarding.
\end{itemize}
\textbf{Parole chiave}: Customer Success, SaaS, Churn reduction, QBR, Onboarding, Upselling}

\cventry{mar 2018–dic 2020}
        {Customer Success Manager}
        {SaaSNow S.r.l.}
        {\url{https://www.saasnow.it}}
        {}
        {\begin{itemize}
\item Responsabile di un portafoglio di 80 clienti SMB nel settore software gestionale, con focus su adozione e soddisfazione.
\item Progettazione di un percorso di onboarding strutturato che ha ridotto il time-to-value medio da 60 a 35 giorni.
\item Monitoraggio di KPI di salute del cliente tramite CRM e dashboard, con interventi proattivi sui clienti a rischio.
\item Gestione delle escalation e coordinamento tra supporto tecnico, prodotto e sales.
\item Contributo alla creazione del team Customer Success e alla definizione dei processi interni.
\end{itemize}
\textbf{Parole chiave}: Onboarding, Time-to-value, KPI, CRM, Escalation}

\cventry{set 2014–feb 2018}
        {Account Manager}
        {Digitalia Consulting}
        {\url{https://www.digitaliaconsulting.it}}
        {}
        {\begin{itemize}
\item Gestione della relazione commerciale con clienti nei settori manifatturiero e servizi, per un portafoglio da 1,2 milioni di euro.
\item Analisi dei bisogni, presentazione di soluzioni digitali e negoziazione dei contratti.
\item Coordinamento di progetti di trasformazione digitale con team interni e partner esterni.
\item Raggiungimento del 115\% degli obiettivi di vendita annuali nel 2017.
\end{itemize}
\textbf{Parole chiave}: Account Management, Negotiation, Digital Transformation, Sales}
\section{Lingue}

\cvline{Italiano}{Madrelingua o bilingue \hfill \textbf{Parole chiave}: Madrelingua}
\cvline{Inglese}{Competenza professionale completa \hfill \textbf{Parole chiave}: C1, Business English}
\cvline{Francese}{Competenza professionale limitata \hfill \textbf{Parole chiave}: A2, Conversazione di base}
\section{Competenze}

\cvline{Customer Success Management}{Esperto \hfill \textbf{Parole chiave}: Onboarding, Fidelizzazione, Customer Journey, QBR, Churn prevention, Upselling}
\cvline{CRM e Analisi Dati}{Avanzato \hfill \textbf{Parole chiave}: Salesforce, HubSpot, Gainsight, Power BI, Tableau, Excel avanzato}
\cvline{Comunicazione e Negoziazione}{Avanzato \hfill \textbf{Parole chiave}: Presentazioni executive, Gestione escalation, Mediazione, Storytelling}
\cvline{Project Management}{Avanzato \hfill \textbf{Parole chiave}: Agile, Jira, Asana, Gestione stakeholder}
\section{Premi}

\cventry{nov 2023}
        {NexaCloud Italia}
        {Customer Success Leader of the Year}
        {}
        {}
        {Riconoscimento interno per i risultati ottenuti nella riduzione del churn e nella costruzione di un modello di customer success scalabile per il segmento enterprise.}

\cventry{feb 2020}
        {SaaSNow S.r.l.}
        {Excellence in Customer Retention}
        {}
        {}
        {Premio assegnato per il miglior tasso di rinnovo del portafoglio SMB e per l'implementazione di un programma di adozione proattiva.}
\section{Certificazioni}

\cventry{mar 2022}
        {SuccessCOACHING}
        {Certified Customer Success Manager (CCSM)}
        {\url{https://www.successcoaching.com/certified-customer-success-manager}}
        {}
        {}

\cventry{giu 2021}
        {Salesforce}
        {Salesforce Certified Administrator}
        {\url{https://trailhead.salesforce.com/credentials/administrator}}
        {}
        {}

\cventry{set 2020}
        {HubSpot Academy}
        {HubSpot Sales Software Certification}
        {\url{https://academy.hubspot.com/courses/sales-software}}
        {}
        {}
\section{Pubblicazioni}

\cventry{giu 2023}
        {Ridurre il churn con l'adozione proattiva}
        {Customer Success Italia}
        {\url{https://www.customersuccessitalia.it/whitepaper/churn-adozione-proattiva}}
        {}
        {\begin{itemize}
\item Analisi delle principali cause di abbandono nei SaaS B2B e delle metriche di early warning.
\item Modello operativo per segmentare i clienti in base al rischio e attivare interventi mirati.
\item Casi studio su onboarding, health score e business review.
\end{itemize}}
\section{Progetti}

\cventry{gen 2022–dic 2022}
        {Percorso di onboarding standardizzato per clienti enterprise e mid-market.}
        {Programma di onboarding scale-up}
        {\url{https://www.nexacloud.it/customer-success}}
        {}
        {\begin{itemize}
\item Ridisegno del percorso di onboarding con template, milestone e responsabilità chiare.
\item Riduzione del time-to-first-value del 30\% e aumento del CSAT post-onboarding da 4,2 a 4,7.
\item Creazione di una libreria di contenuti self-service e formazione per i clienti.
\end{itemize}
\textbf{Parole chiave}: Onboarding, Time-to-value, CSAT, Self-service}

\cventry{mar 2023–set 2023}
        {Modello predittivo per identificare i clienti a rischio di abbandono.}
        {Health score per il portafoglio clienti}
        {\url{https://www.nexacloud.it/insights/health-score}}
        {}
        {\begin{itemize}
\item Definizione di indicatori di adozione, engagement e soddisfazione ponderati per segmento.
\item Integrazione del modello con CRM e dashboard operative per i Customer Success Manager.
\item Aumento del 25\% degli interventi proattivi sui clienti a rischio.
\end{itemize}
\textbf{Parole chiave}: Health score, Analytics, Churn, CRM}
\section{Interessi}

\cvline{Innovazione digitale}{AI generativa, Product-led growth, Design thinking}
\cvline{Sport e benessere}{Running, Yoga, Trekking}
\cvline{Comunità e formazione}{Mentoring, Public speaking}
\section{Volontariato}

\cventry{set 2021–Presente}
        {Mentor}
        {Women in Tech Italia}
        {\url{https://www.womenintechitalia.it}}
        {}
        {\begin{itemize}
\item Mentoring individuale per giovani professioniste che desiderano intraprendere carriere nel customer success e nelle vendite SaaS.
\item Partecipazione a incontri di orientamento e workshop su comunicazione, negoziazione e crescita professionale.
\item Supporto alla community nella creazione di percorsi di sviluppo e networking.
\end{itemize}}
    
\end{document}